\section{Related Work}
\label{sec:related_work}

\subsection{Alignment and Preference Optimization}
Optimizing Large Language Models (LLMs) to align with human intent has predominantly relied on Reinforcement Learning from Human Feedback (RLHF) \citep{ouyang2022training}, utilizing Proximal Policy Optimization (PPO) \citep{schulman2017proximal}. Recently, Direct Preference Optimization (DPO) \citep{rafailov2023direct} emerged as a stable alternative that mathematically bypasses the need for an explicit reward model. Building on this, recent advancements like RRHF \citep{yuan2023rrhf}, ORPO \citep{hong2024orpo}, and SimPO \citep{meng2024simpo} have further simplified alignment via reference-free objectives. Recent work adds adaptive batch-wise scheduling \citep{zhang2025adaptive} and constrained-alignment methods \citep{chen2025constrained}, addressing algorithmic biases in RLHF \citep{ji2024limits, shrivastava2024algorithmic}. A parallel line \emph{unifies} SFT and RL into a single objective \citep{liu2025uft, wang2024uft}; our contribution is orthogonal and composable with it, keeping the standard SFT$\to$DPO pipeline and changing only how preference pairs are labelled. Like RLAIF \citep{lee2023rlaif} our labels need no human annotation; unlike it, they come from an NLI model or a pairwise reasoning judge rather than a general-purpose LLM preference.
\subsection{Reasoning Quality and Faithfulness Rewards}
Rewarding the logical quality of a generation, rather than human preference, has been studied with process-based feedback \citep{uesato2023verify}, checks on decision pivots in chain-of-thought \citep{wang2025pivots}, and entailment rewards for summarisation faithfulness \citep{pasunuru2018multireward, roit2023factually}. Retrieval \citep{lewis2020retrieval} and hallucination detection \citep{li2023haludetect} target factual errors; we do not address hallucination beyond reporting how often SFT outputs contain fabricated URLs. Our concern is narrower: whether an entailment reward measures reasoning or can be met by restating the answer, a risk also noted for metric-driven evaluation generally \citep{schaeffer2023mirage}.

\subsection{LLM-as-a-Judge and Evaluation Bias}
The use of powerful instruction-tuned models like GPT-4 as automated evaluators \citep{zheng2024judging} is widespread but prone to significant biases, including verbosity bias \citep{wang2023large} and self-preference bias \citep{pan2024evalbias}. The verbosity-bias confound has been formally analysed in DPO itself: \citep{park2024lengthbias} show that standard DPO and its $\beta$-tuned variants \citep{park2024betadpo} systematically conflate explanation length with quality, and propose length-aware regularisation as a partial remedy. We avoid LLM-judge labels for training the NLI-reward policies, and where we use LLM judges for evaluation we compare outputs of matched length and use several judges from different model families. Table~\ref{tab:pairwise_1} shows that, without such controls, a GPT-4o-mini judge picks the longer of two outputs in every comparison we ran.
